\subsection{معرفی و انگیزه}
این مکانیزم بر مبنای مبانی نظری پیشرفته در ادبیات توجه و فشرده‌سازی \cite{eckart1936approximation} توسعه یافته است.

افزایش مقیاس مدل‌های زبانی بزرگ، چالش‌های بنیادینی را در حافظه فعال و توان عملیاتی استنتاج پدید آورده است. بخش عمده‌ای از زمان تاخیر محاسباتی ناشی از گذر متوالی تانسورهای پنهان از میان ده‌ها لایه ترنسفورمر یکسان است. بررسی همبستگی میان‌لایه‌ای نشان می‌دهد که بازنمایی لایه‌های مجاور در اعماق میانی ترنسفورمرها دارای افزونگی معنایی بالایی بوده و شباهت کسینوسی حالت‌های پنهان میان لایه‌های متوالی اغلب به بیش از ۰٫۹۵ می‌رسد.

ماژول \lr{come\_compressor} (مخفف \lr{CoMe: Layer Concatenation and Shedding with Learnable Inter-Layer Interpolation Adapters}) با هدف هرس ساختاری لایه‌ها (\lr{Structured Layer Shedding}) به همراه پل‌های درون‌یابی آموزش‌پذیر میان‌لایه‌ای (\lr{Learnable Inter-Layer Interpolation Bridges}) طراحی گردیده است. رویکرد مرسوم حذف لایه‌ها به دلیل شکست پیوستگی نمایش و شوک ناگهانی به جریان پسماند (\lr{Residual Stream})، افت دقت شدیدی را تجربه می‌کند. ماژول \lr{come\_compressor} با جایگزینی شکاف‌های ناشی از حذف لایه‌ها توسط پل‌های کم‌مرتبه سبک با مقداردهی خنثی (\lr{Identity-Initialized Bridges})، پیوستگی جریان نمایش پنهان را حفظ کرده و اجازه می‌دهد مدل فشرده‌شده بدون نیاز به پیش‌آموزش مجدد، به سرعت همگرا شود.
